% Part: first-order-logic
% Chapter: beyond
% Section: other-logics

\documentclass[../../../include/open-logic-section]{subfiles}

\begin{document}

\olfileid{fol}{byd}{oth}

\olsection{Nog meer soorten logica}

Een nieuwe logica maken is niet zo moeilijk. Je kiest een eigen syntaxis, maakt
een systeem met afleidingsregels en geeft aan wat de formules betekenen. Dat
laatste is de semantiek. Je hebt wel wat vindingrijkheid nodig als je wilt dat
het !!{derivation} systeem volledig is voor die semantiek. Ook kan het moeite
kosten om anderen ervan te overtuigen dat je logica echt interessant is. Maar
je kunt er uren plezier aan beleven om uit te zoeken welke wiskundige en
computationele eigenschappen de logica heeft.

In de afgelopen decennia zijn er heel veel formele logica's bij gekomen. Fuzzy
logic beschrijft redeneringen over vage eigenschappen. Probabilistische logica
beschrijft redeneringen over onzekerheid en kansen. Defaultlogica's en
niet-monotone logica's beschrijven redeneringen waarvan een conclusie later kan
vervallen. Het gaat om ``redelijke'' gevolgtrekkingen die gelden zolang nieuwe
informatie ze niet onderuit haalt. Epistemische logica's gaan over kennis.
Causale logica's gaan over oorzaak en gevolg. ``Deontische'' logica's gaan over
morele en ethische verplichtingen: wat moet, wat mag en wat verboden is. De
eerste reden om zo'n systeem te maken kan filosofisch, wiskundig of
computationeel zijn. Daarom worden zulke systemen op allerlei plaatsen
bestudeerd: binnen de wiskundige logica, de filosofische logica, de kunstmatige
intelligentie, de cognitiewetenschap en daarbuiten.

De lijst houdt niet op en de mogelijkheden lijken eindeloos. Misschien lukt het
nooit om Leibniz' droom waar te maken en al het menselijk redeneren terug te
brengen tot berekeningen---maar we kunnen het wel blijven proberen.

\end{document}
